\documentclass[sigconf]{acmart}
\emergencystretch=2em   

\settopmatter{printacmref=false}   
\settopmatter{printfolios=false}

\title[TestVDB: source-grounded falsification for VDBMS API documentation-implementation consistency]{TestVDB: Source-Grounded Falsification of LLM-Derived Behavioral Claims for Documentation-Implementation Consistency Testing of Vector Databases}

\author{Authors TBD}
\email{tbd@example.org}
\affiliation{\institution{TBD (anonymity depends on venue)}\country{TBD}}

\begin{abstract}
Vector Database Management Systems (VDBMSs) store the embeddings that retrieval-augmented Large Language Model (LLM) applications depend on, yet a large class of their defects lacks a practical oracle. We study \emph{documentation-implementation defects}: cases where a VDBMS silently accepts an input or behavior that violates its API documentation, the natural-language description of intended behavior, such as accepting out-of-range index parameters (\texttt{nprobe=0}, \texttt{ef=0}) or returning success where the documentation prescribes rejection. These defects corrupt query semantics; the large majority produce no crash, so fuzzers miss them.

Of the 107 issues TestVDB surfaced (49 true-positive defects: 15 fixed via merged PR, 16 with open fix-PRs, and 18 maintainer-acknowledged but not yet fixed), about 85\% are documentation-implementation defects that differential, metamorphic, and property-based oracles cannot reach; this is the composition of our findings, not a population estimate. The reason is structural: these accept/reject decisions cannot be mechanically checked, because the documented boundary is natural-language rather than formal. We therefore adopt an LLM as the semantic judge; prior REST-API oracle work avoids this setting by keeping deterministic, executable assertions. The judge's errors split in two. Family-specific self-preference is mitigated by cross-model validation. \emph{Task-intrinsic} errors, where ambiguous documentation makes different LLM families infer the same wrong claim, are not: in an eighteen-clause probe on Milvus and Qdrant, a different family reproduced GLM's over-strict claim in 6, and the over-strict phenomenon concentrates in optional-default APIs (none on Weaviate, whose docs state explicit bounds). The implementation, as the most accessible source of actual behavior, is what resolves them.

We instantiate this as TestVDB, which falsifies LLM-derived behavioral claims against source. TestVDB surfaced 107 candidate issues across three VDBMSs; 49 are true-positive defects (15 merged-PR-fixed, 16 with open fix-PRs, 18 acknowledged but unfixed), and a controlled retrospective on Milvus and Qdrant shows the source-grounded dev-reviewer reaches 67\% precision and 74\% recall (3-run ensemble on 48 adjudicated candidates), versus 37\% recall without the source anchor.
\end{abstract}

\begin{CCSXML}
<ccs2012><concept><concept_id>10011007.10011074.10011099</concept_id><concept_desc>Software and its engineering~Software testing and debugging</concept_desc><concept_significance>500</concept_significance></concept></ccs2012>
\end{CCSXML}
\ccsdesc[500]{Software and its engineering~Software testing and debugging}

\keywords{Vector database, documentation-implementation consistency testing, test oracle, LLM-derived oracle, LLM-as-judge, source-grounded falsification}

\begin{document}

\maketitle

\section{Introduction}

Vector Database Management Systems (VDBMSs) store the embeddings that retrieval-augmented LLM applications depend on. Their defects are costly: an empirical study~\cite{bugstudy25} finds that more than half of VDBMS bugs manifest as functional failures, and the VDBMS testing roadmap~\cite{roadmap25} attributes about 43\% of VDBMS bugs to incorrect behavior and identifies oracle definition as a key challenge. Existing VDBMS fuzzers~\cite{vdbfuzz26} apply classical fuzzing~\cite{manes21} but detect only crashes. We target a prevalent subset: \emph{documentation-implementation defects}, where a VDBMS silently accepts an input or produces a behavior that violates its API documentation, such as accepting \texttt{nprobe=0} or out-of-range HNSW index parameters such as \texttt{ef=0}, or returning success where the documentation prescribes rejection. These defects corrupt query semantics and can return data the documentation intended to exclude (for example, a negative score threshold that disables a filter and returns all matches). Most produce no crash (the large majority of the 49 true-positive defects in our data).

The oracle problem~\cite{barr15} is acute here. Across 107 submitted VDBMS issues (49 true-positive defects), about 85\% are, by our classification, documentation-implementation defects that the standard oracle families cannot reach (Table~\ref{tab:exclusion}). Differential testing catches mathematical-invariant violations across vendors but cannot adjudicate accept/reject, which diverges across VDBMSs by design and leaves no cross-vendor reference; the roadmap flags differential testing as challenging in this setting~\cite{roadmap25}. Metamorphic relations~\cite{metmap24} address result correctness but not input-acceptance decisions, because a metamorphic relation is an output relation, not an input-acceptance transform. Property-based testing needs a machine-checkable property and an OpenAPI schema; VDBMS endpoints serve no schema that encodes these documentation constraints, and several serve no OpenAPI at all. The root cause is that many of these accept/reject decisions cannot be mechanically checked, because the documented boundary is natural-language and ambiguous rather than formal.

\begin{table*}[t]
\caption{Oracle candidates by the defect class they reach, and why the documentation-implementation residual leaves an LLM as the practical oracle. Rows 1 to 4 use a deterministic oracle; row 5 (REST doc/spec-derived assertion oracles) also keeps a deterministic oracle and so never enters this regime; only row 6 does.}
\label{tab:exclusion}
\small
\begin{tabular}{@{}p{0.22\textwidth}p{0.30\textwidth}p{0.40\textwidth}@{}}
\toprule
Candidate oracle & Reaches (defect class) & Why it misses the documentation-implementation residual \\
\midrule
Crash (VDBFuzz~\cite{vdbfuzz26}) & crash/hang & most documentation-implementation defects do not crash (the large majority of the 50); silent-accept defects are invisible to a crash oracle (controlled head-to-head, \S\ref{sec:eval}) \\
Differential testing & math invariants across vendors (our model-free subclass, \S\ref{sec:eval}) & cross-vendor accept/reject diverges by design; no reference~\cite{roadmap25} \\
Metamorphic relations~\cite{metmap24} & result correctness (top-$k$ monotonicity, recall vs.\ $ef$) & an MR is an output relation; documentation consistency is an input accept/reject decision, and no transform preserves it \\
Property-based testing~\cite{claessen00,schemathesis,quickrest20} & math properties + schema conformance & needs a machine-checkable property and an OpenAPI schema; VDBMS endpoints serve no schema that encodes these constraints, and several serve no OpenAPI \\
REST doc/spec-derived oracles (AGORA+~\cite{agoraplus25}, SATORI~\cite{satori25}, MASTOR~\cite{mastor26}) & status/field/invariant assertions from OpenAPI, traces, or source & reliable extraction from low-ambiguity structured sources; no falsification needed \\
\textbf{LLM-derived oracle (TestVDB)} & \textbf{accept/reject vs.\ API documentation} & \textbf{many documentation semantics are not mechanically checkable, leaving an LLM as the practical oracle for the residual; the reliability problem this introduces is the subject of \S\ref{sec:regime}} \\
\bottomrule
\end{tabular}
\end{table*}

The absence of a mechanical oracle for these cases places VDBMS documentation-implementation consistency in a setting where the LLM serves as the semantic oracle (Section~\ref{sec:regime}): an LLM must both extract behavioral claims from the documentation and, where extraction is inapplicable, judge documentation consistency directly. The gap is in source ambiguity: prior REST-API oracle work~\cite{agoraplus25,satori25,mastor26} extracts from structured sources where constraints are explicit, while VDBMS documentation is ambiguous and LLM interpretation may produce wrong claims. The errors this introduces split into a family-specific layer that cross-model validation mitigates and a \emph{task-intrinsic} layer it cannot, because the ambiguity is in the documentation. \emph{Source-grounded falsification} (Section~\ref{sec:design}) resolves the task-intrinsic layer by treating these claims as refutable hypotheses and checking them against the implementation's actual behavior. We instantiate it in TestVDB (Section~\ref{sec:design}), a detector that extracts behavioral claims from documentation, probes VDBMS endpoints, and falsifies LLM verdicts against source.

\paragraph{Contributions.}
\begin{itemize}
  \item \textbf{The role of the LLM (\S\ref{sec:regime}).} We show that VDBMS documentation, unlike structured specifications, forces an LLM into both extraction and judgment roles, that the resulting behavioral claims are unreliable in two distinct ways (family-specific and task-intrinsic), and that this source-ambiguity gap (low-ambiguity sources yield reliable assertions, whereas ambiguous natural-language documentation yields claims that may be wrong) is what separates VDBMS documentation-implementation consistency testing from prior REST-API oracle work.
  \item \textbf{Quantification of the documentation-implementation residual.} About 85\% of the documentation-implementation defects we found are unreachable by differential, metamorphic, or property-based oracles, mapping where each classical class fails.
  \item \textbf{Task-intrinsic documentation-interpretation errors and the source-grounded counter.} We separate family-specific from task-intrinsic documentation-interpretation errors and, on an eighteen-clause probe on Milvus and Qdrant, find that cross-model judging misses 3 of the 6 task-intrinsic clauses while source-grounded falsification catches all 18, and that over-strict concentrates in optional-default APIs (Milvus and Qdrant's search parameters) and is absent where documentation states explicit bounds (Weaviate, whose doc-code gaps are documentation-implementation bugs instead).
  \item \textbf{TestVDB and scale.} TestVDB surfaced 107 candidate issues across three VDBMSs; 49 are true-positive defects (15 merged-PR-fixed, 16 with open fix-PRs, 18 acknowledged but unfixed). A controlled retrospective on Milvus and Qdrant shows the source-grounded dev-reviewer reaches 67\% precision and 74\% recall (3-run ensemble).
\end{itemize}

\section{Background and Problem Setup}

Barr et al.~\cite{barr15} classify test oracles by the artifact that supplies the pass/fail verdict: \emph{specified} (from a specification or contract), \emph{derived} (from execution, as in differential, metamorphic, or invariant testing), \emph{implicit} (a crash), or \emph{none} (a human). TestVDB uses two: a documentation-derived oracle for documentation-implementation defects, in the specified class\footnote{VDBMS API documentation, unlike the formal assertions of Design by Contract~\cite{meyer92dbc}, is natural-language prose; this informality is the source of both the need for an LLM and the reliability problem we address.}, and a model-free invariant oracle for mathematical-bound violations, in the derivable class. The subset we care about is documentation-implementation consistency, and the central observation is that on this subset the first three classes fail, leaving only the human oracle, or its scalable proxy, an LLM.

Conformance testing originated in protocol standards~\cite{iso9646}. We study a documentation-implementation consistency problem in VDBMSs: cases where the API silently accepts an input or behavior that diverges from its API documentation. We separate documentation-implementation consistency from correctness. \emph{Documentation-implementation consistency} asks whether the API's accept/reject behavior matches its documentation. \emph{Correctness} asks whether a returned result is mathematically right, such as approximate nearest neighbor (ANN) recall or ranking. TestVDB targets documentation-implementation consistency; result correctness of vector search remains open and is not our claim~\cite{roadmap25}.

For a substantial portion of these cases, the documented boundary is natural-language prose rather than a formal grammar, so the verdict cannot be mechanically checked without semantic interpretation. We formalize the consequence in Section~\ref{sec:regime}.

\section{The Role of the LLM and Its Reliability Problem} \label{sec:regime}

VDBMS API documentation, unlike the structured sources that REST-API oracle tools rely on (OpenAPI specifications, execution traces, or source code), is natural-language prose. This forces an LLM into the testing pipeline~\cite{hou23llmse} in two ways. First, the LLM must extract behavioral claims from the documentation, translating vague descriptions such as ``optional, default 1'' into checkable clauses such as ``parameter $\geq 1$.'' Second, when the documentation leaves the treatment of some inputs entirely unstated, the LLM must judge documentation consistency directly, because no extractable clause exists. In both roles, the LLM's output is a behavioral claim that may be wrong: the extraction may over-formalize the documented intent, and the direct judgment may hallucinate~\cite{ji23hall}. A human can perform these roles but does not scale, so in practice both extraction and judgment fall to a language model. This setting combines extraction from ambiguous documentation with, where extraction fails, direct semantic judgment; in both roles, the oracle is LLM-derived and therefore unreliable.

The contrast with recent REST-API oracle work is about source ambiguity, not extraction mechanism. AGORA+~\cite{agoraplus25}, SATORI~\cite{satori25}, and MASTOR~\cite{mastor26} extract oracles from structured sources (OpenAPI specifications, execution traces, or source code) where the constraints are explicit and the LLM's task is to transcribe them into assertions. Even though SATORI and MASTOR also use LLMs, the low ambiguity of their sources means the extracted assertions are reported as reliable, with low false-positive rates. TestVDB extracts from natural-language documentation, where the constraints are implicit and ambiguous and the LLM's task is to interpret intent. Interpretation from ambiguous text is unreliable: the LLM may stably produce a wrong claim: the same wrong claim across runs and across model families, because the ambiguity is in the shared input, not the model. This stable misinterpretation is the gap that source-grounded falsification addresses. MASTOR is the closest in using source, but it tests what the implementation does, with source as the reference, and so cannot detect a gap between the documentation and the code; TestVDB tests what the documentation prescribes, with source as the reference for actual behavior, and that gap is exactly what it targets.

Entering this setting exposes a reliability problem with two layers. The first is family-specific: when one LLM family both extracts behavioral claims from the documentation and judges documentation consistency, the two roles share biases, and the judge tends to confirm the extractor's errors. This is the LLM-as-judge self-preference phenomenon (Panickssery et al.~\cite{panickssery24}), and cross-model validation, in which a second family judges, mitigates it. The second layer is deeper. When the documentation itself is ambiguous, different LLM families extract the same wrong claims. Two failure modes result. In \emph{over-strict} extraction, a parameter documented as ``optional, default 1'' is over-formalized by both families as ``must be at least 1,'' when the real semantics is that the value 0 selects the default. In \emph{fabricated} extraction, a parameter whose documentation illustrates a power-of-two value is hallucinated by both families as ``must be a power of two,'' a constraint the documentation never states. In both cases the error originates in the shared input rather than in either model, so cross-model validation cannot remove it.

We scope the stability claim precisely. Task-intrinsic stability is an extraction-level property defined across model families: a clause is task-intrinsic when a second family's independent formalization of the same documentation is also over-strict on the same parameter. This is distinct from the intra-judge self-inconsistency that Haldar et al.~\cite{haldar25} report for LLM-as-a-judge ratings, where a single judge's verdict on the same input varies across runs---a sampling-noise phenomenon in the judgment step. Our claim concerns extraction: when the documentation is ambiguous, different families converge on the same wrong clause, and that convergence is what cross-model validation cannot break.

We confirmed the split with two probes on eighteen over-strict clauses (thirteen from Milvus, five from Qdrant v1.18.2, live-probed to confirm over-strict). We first fed each raw documentation passage to two LLM families, GLM and DeepSeek, and asked each to formalize the documented behavior independently; DeepSeek reproduced an over-strict clause on the same parameter in 6 of the 18 (2 of 13 Milvus, 4 of 5 Qdrant), the task-intrinsic subset; a clause counts as task-intrinsic when the second family's independent formalization is also over-strict on the same parameter, a comparison at the parameter level rather than verbatim. We then tested the operational alternative directly: we asked DeepSeek to judge each of GLM's over-strict clauses against the documentation. Cross-model judging caught 8 of the 18 but missed 3 of the 6 task-intrinsic ones, where DeepSeek read the ambiguous documentation as GLM had. Source-grounded falsification, by contrast, contradicts every over-strict clause, because the implementation accepts the value the clause rejects. The original probe was a twelve-clause pilot; two scaling passes extend it to eighteen (Section~\ref{sec:eval}).

The implication is the basis for our method. Cross-model validation covers the family-specific subset but not the task-intrinsic one. We also tried multi-perspective judging---four specialized LLM judges (documentation, evidence, severity, novelty) with a vote---but on the retrospective (\S\ref{sec:eval}) it reaches only 15\% recall at 80\% precision: the judges share the same documentation ambiguity and converge on the same wrong claims, so voting cannot break task-intrinsic agreement. The task-intrinsic subset requires a source of actual behavior that does not depend on the documentation; the most accessible one is the implementation. This motivates source-grounded falsification, developed in Section~\ref{sec:design}: treat the LLM-derived behavioral claims as refutable hypotheses and falsify them against source, where the documentation expresses the intended behavior and the source expresses the actual behavior.

\section{TestVDB Design} \label{sec:design}

TestVDB instantiates source-grounded falsification. It treats LLM-derived behavioral claims as a set of clauses, each a hypothesis that the API enforces some documented constraint, and uses the implementation's actual behavior to falsify them.

The falsification rule is concrete. For a clause that asserts the API rejects inputs where a parameter violates a constraint (for example, \texttt{shardsNum} $\geq 1$), the dev-reviewer (the pipeline's source-reading agent) examines the implementation. If the source shows that \texttt{shardsNum = 0} selects the default shard number, the clause is over-strict and the verdict that relied on it is falsified, a false positive suppressed. If the source shows no such intended semantics yet the implementation still accepts a value the documentation prescribes rejecting, the violation stands as a real defect.

This use of source is the opposite of MASTOR's~\cite{mastor26}: both are source-grounded, but MASTOR reads source to generate oracles that encode implemented behavior and so cannot detect a gap between the documentation and the code, whereas TestVDB reads source to falsify documentation-derived clauses and targets exactly that gap. The novelty over MASTOR is therefore the asymmetric direction in which source is used---as an independent falsifier of documentation-derived claims rather than as the oracle itself---not source grounding per se (Section~\ref{sec:related}).

The pipeline has five stages. An LLM first extracts behavioral claims from the documentation as clauses. Attack agents then generate boundary inputs targeting each clause. An LLM-derived oracle judges whether each observed response conforms to the documentation; this is the step whose two-layer unreliability we described in Section~\ref{sec:regime}. The dev-reviewer then falsifies the clauses and the LLM's verdicts against source. A novelty gate finally removes duplicates and known issues.

The dev-reviewer is a source-grounded falsifier: it acts as a VDBMS maintainer would---reproduce each candidate, cross-check it against source, and try to disprove it. For each candidate it applies three checks: (1) \emph{independently reproducible}---does the candidate reproduce under a clean probe, free of tooling artifacts? (2) \emph{evidence sufficient}---do the response, the VDBMS state, and the source code jointly establish the documented-implementation gap? (3) \emph{falsifiable}---can the candidate be disproven by an alternative documented interpretation, such as a by-design behavior or a wont-fix pattern? A candidate that survives all three is adjudicated a defect; one that fails any check is suppressed. The checks draw on three anchors: a clean reproduction that re-probes the live API, a source-grounded anchor that is the primary falsifier, and an optional threat-model cross-check that compares each candidate against maintainer-acknowledged by-design and wont-fix patterns. This separation matters because the same LLM family produces both the behavioral claims and the verdicts; the implementation, an independent information source, is what falsifies them.

\section{Implementation} \label{sec:impl}

TestVDB runs as a multi-agent pipeline~\cite{he2025lma} on the Claude Code runtime\footnote{\url{https://docs.anthropic.com/en/docs/claude-code}, accessed 2026-07.}, with each of the 20 agents defined by a task-structured role prompt and dispatched to a GLM-5.2 backbone\footnote{GLM Technical Report; see \url{https://www.z.ai/blog}, accessed 2026-07.} (via the bigmodel Anthropic-compatible API) under the Claude Code runtime's default sampling, with no fixed random seed; we have not measured run-to-run variance and flag this as a limitation. Target VDBMSs run in pinned Docker versions (Milvus 2.6.19 for the single-LLM and single-layer ablations; the full version matrix is in the artifact). A complete run per target exercises the 20 agents over tens of generated candidates; wall-clock is dominated by the dev-reviewer's source-grounding step, repository clone and source retrieval, and live Docker re-probes, rather than by raw LLM latency. The 107-submission study plus the evaluation batches amount to on the order of $10^4$ LLM calls and roughly $\$10$ per target at current API pricing\footnote{Based on public provider pricing as of 2026-07; the manual-testing comparison reflects the authors' experience.}, comparable to a few hours of manual boundary testing. The full prompts, target versions, and per-token accounting are in the artifact, which we will release at a persistent URL upon acceptance.

\section{Evaluation} \label{sec:eval}

We evaluate TestVDB on three VDBMSs (Milvus, Qdrant, Weaviate) and ask three questions. \textbf{RQ1:} how many documentation-implementation defects does TestVDB surface, and how large is the residual beyond classical oracles? \textbf{RQ2:} does source-grounded falsification suppress false positives? \textbf{RQ3:} can cross-model validation resolve task-intrinsic documentation-interpretation errors, or is source needed?

\paragraph{RQ1: yield and the documentation-implementation residual.}
TestVDB surfaced 107 candidate issues across the three systems; 49 are true-positive defects (15 fixed via merged PR, 16 with open fix-PRs, and 18 maintainer-acknowledged but not yet fixed: 8 still open with an accepted label and 10 acknowledged but stale-closed without a merged PR). The remaining 58 comprise 23 by-design/rejected and 35 still pending maintainer adjudication. Submissions span Milvus (51), Weaviate (30), and Qdrant (26); true positives span Milvus (22), Weaviate (13), and Qdrant (14) (Table~\ref{tab:yield}). The 10 stale-closed TPs carry an accepted label or drew no maintainer denial; we count them as true defects because the maintainer acknowledged the issue rather than rejecting it, even though no merged-PR fix exists. To quantify the residual, we classified each submitted issue by fault model: classical-addressable (a mathematical-invariant or crash that a differential, metamorphic, or crash oracle could in principle find), documentation-implementation (an accept/reject violation only a semantic judge reaches), or concurrency (a race or atomicity violation)\footnote{The per-issue mapping---all 107 submissions assigned to one of these three fault models with a one-line rationale, plus a maintained count of each---is in the artifact, so the 85\% composition is auditable rather than asserted.}. About 85\% of the issues we submitted are, by this classification, documentation-implementation defects that classical oracles cannot reach (89\% on the 49-TP subset). Of the 72 maintainer-adjudicated submissions (49 true-positive, 23 by-design/rejected), the yield precision is 68.1\% (Wilson 95\% CI $[56.6\%, 77.7\%]$); treating all 35 still-pending submissions as false positives gives a worst-case bound of 45.8\% (Wilson 95\% CI $[36.7\%, 55.2\%]$). The remainder split between classical-addressable cases (about 10\%) and concurrency defects (about 5\%). A full-space estimation with capture-recapture or an unbiased defect sample is future work. As a structural check, we ran a classical-oracle suite on Qdrant v1.18.2: a metamorphic relation set (distance symmetry, top-$k$ monotonicity, the \texttt{COSINE} bound, self-similarity) found no violations on this version and, by construction, no documentation-implementation defects, since metamorphic relations relate outputs rather than input-acceptance decisions. The mathematical-invariant violations such checks surface---9 of the 107 submissions, including \texttt{COSINE} distance above 1 for identical vectors and an index returning 2 of 25 matching points---are detected by a model-free invariant subclass that probes hard mathematical bounds directly, with no documentation extraction or LLM judgment; the documentation-implementation residual lies outside classical reach. A controlled head-to-head with VDBFuzz~\cite{vdbfuzz26} sharpens this residual beyond the disjoint-classes observation. On Qdrant v1.18.2, we ran VDBFuzz (default configuration), which executed over 26{,}000 mutated requests and found 0 crashes; the two Qdrant case-study crashes are fixed in this version, so the result reflects version state rather than an insufficient budget. We therefore ran a bidirectional reachability probe on versions where each tool's target defect is reproducible (Table~\ref{tab:versions}; v1.4.0 reproduces the integer-overflow crash; v1.18.0 predates the May~2026 fix that resolves \#9045; v1.18.2 is the version TestVDB targeted, in which both are fixed). On the older Qdrant v1.4.0, where VDBFuzz's integer-overflow crash on a vector \texttt{size} of $2^{63}$ is still live, TestVDB's boundary probe flags it as a documentation-implementation violation: the OpenAPI declares \texttt{size}~$\geq 1$ with no maximum, so the value is documented-valid, yet the implementation panics. On Qdrant v1.18.0, where TestVDB's acknowledged empty-vector defect (\#9045)\footnote{\url{https://github.com/qdrant/qdrant/issues/9045}} is live---\texttt{wait=false} silently accepts a zero-length vector---VDBFuzz's empty-vector template reports no failure: the template exercises \texttt{wait=true} (hardcoded in every Qdrant points-upsert request in VDBFuzz's released template suite\footnote{Every points-upsert template in VDBFuzz's Qdrant suite~\cite{vdbfuzz26} uses \texttt{?wait=true}; the \#9045 violation lives on the \texttt{wait=false} path, which no template probes.}), the path that correctly rejects, and does not probe \texttt{wait=false} where the violation lives; even on that path the response is a silent HTTP 200, which a crash oracle cannot flag. We therefore read the reverse direction as a limitation of VDBFuzz's current template and input coverage, not as a fundamental property of crash oracles. Each direction is $n{=}1$; we treat these as hypothesis-generating controlled cases rather than a generalized result, and they suggest that a documentation-implementation oracle can reach a crash-class defect whose input violates a documented bound, while a crash oracle, under VDBFuzz's current templates, does not reach a documentation-implementation defect that manifests as a silent accept; the structural hypothesis is in Section~\ref{sec:disc}.

\begin{table}[t]
\caption{Bidirectional reachability: the Qdrant version on which each tool's target defect reproduces, and the fix state. Each direction is $n{=}1$.}
\label{tab:versions}
\small
\begin{tabular}{@{}p{0.13\columnwidth}p{0.52\columnwidth}p{0.30\columnwidth}@{}}
\toprule
Version & Reproduces & Fix state \\
\midrule
v1.4.0 & VDBFuzz's integer-overflow crash (\texttt{size}=$2^{63}$) & fixed in v1.5.0 \\
v1.18.0 & TestVDB's \#9045 (\texttt{wait=false} accepts a zero-length vector); \#7967 panic under load & fixed May~2026 \\
v1.18.2 & --- (TestVDB's target version) & both fixed \\
\bottomrule
\end{tabular}
\end{table}

\begin{table}[t]
\caption{Issues submitted and maintainer-acknowledged per VDBMS.}
\label{tab:yield}
\small
\begin{tabular}{@{}lrr@{}}
\toprule
VDBMS & Submitted & Acknowledged \\
\midrule
Milvus & 51 & 22 \\
Weaviate & 30 & 13 \\
Qdrant & 26 & 14 \\
\midrule
Total & 107 & 49 \\
\bottomrule
\end{tabular}
\end{table}

\paragraph{RQ2: effect of source-grounded falsification.}
On a controlled retrospective over 48 maintainer-adjudicated candidates (27 true-positive, 21 by-design/rejected; Milvus 32, Qdrant 16), the contract-grounded dev-reviewer with a 3-run any-confirmed ensemble reaches 65\% accuracy, 67\% precision (Wilson 95\% CI $[49\%, 81\%]$), and 74\% recall, against a single-LLM baseline (no source-grounded anchor) of 48\%/56\%/37\%; the source-grounded contract anchor is what drives the recall gain ($37\%\to74\%$). Single runs vary widely---across five independent runs, recall ranges $15$--$78\%$ (accuracy $44$--$65\%$, precision $50$--$73\%$), because some runs are conservative and confirm few candidates rather than admitting many false positives. We report the any-confirmed (union) ensemble as the operating point: the dev-reviewer is a falsifier, so a candidate that survives any independent falsification attempt is more likely a true defect, and under-confirmation is costlier than forwarding a false positive for human triage; the union restores recall ($74\%$ at 3 runs; $85\%$ at 5 runs at $62\%$ precision). Majority voting ($\ge 3$ of 5) reaches $64\%$ precision and $26\%$ recall---precision comparable to the single-run band ($50$--$73\%$) but recall too low to be useful, because conservative runs dominate the vote. The union rule, by surfacing candidates that any run confirms, captures signals from high-recall runs without inflating false positives beyond the per-run precision band. Per vendor, Milvus reaches $69\%/73\%/80\%$ and Qdrant $56\%/50\%/57\%$\footnote{The 48 candidates, all five dev-reviewer runs, and the corrected ground truth (27 TP + 21 FP, after the Qdrant reclassification) are in the artifact, together with a stability analysis reporting single-run variance and a five-run ensemble; the earlier single-run, single-vendor figures (81\% false-positive suppression, 69.2\% precision, 96.7\% recall) rested on a 16-candidate control whose data is no longer recoverable and are superseded by this reproducible cross-vendor ensemble.}. The anchors contribute at different points. On a Milvus re-probe of 27 candidates, the clean-reproduction anchor alone kills tooling-artifact false positives at 12/14 but by-design ones at only 2/11. A three-condition ablation on a 12-FP/4-TP control (Milvus v2.6.19, dev-reviewer-v2, blind) then isolates the dev-reviewer's two anchors: source alone suppresses 9/12 FPs (75\%), the threat-model anchor alone 6/12 (50\%, unstable on state/concurrency FPs), and their union 11/12 (91\%), each retaining all 4 TPs, confirming the source-grounded anchor as the dominant contributor. The full per-anchor breakdown is in the artifact. As a cross-model check, DeepSeek re-running the dev-reviewer on twenty candidates across three studies---six from the retrospective plus fourteen newly reviewed blind to GLM-5.2's rationale; the twenty form a diversity-stratified, non-random sample chosen to span input-validation, upsert-semantics, idempotent-drop, correct-reject, and dynamic-field subtypes---agrees with GLM-5.2 on all twenty (Cohen's $\kappa = 1.0$), so the source-grounded verdict does not appear family-specific when source evidence is explicit. RQ3 next isolates the subset of these errors that cross-model validation cannot reach, where source is the only resolver.

\paragraph{RQ3: cross-model validation vs source on task-intrinsic errors.}
The central claim is that source-grounded falsification resolves task-intrinsic documentation-interpretation errors that cross-model validation cannot. We tested this on eighteen over-strict clauses (Table~\ref{tab:e2}): the original twelve (nine from GLM's Milvus collection-creation formalization plus three from Qdrant v1.18.2, confirmed over-strict by live probe---\texttt{timeout=0}, \texttt{group\_size=0}, and an out-of-range \texttt{score\_threshold} are accepted, while \texttt{limit=0}, \texttt{shard\_number=0}, and \texttt{replication\_factor=0} are rejected) plus six added in two scaling passes---Milvus \texttt{ef}, \texttt{nprobe}, \texttt{level}, and \texttt{replicaNumber}, and Qdrant \texttt{m} and quantization \texttt{bits}---each formalized by the same contract-formalizer and live-probe-confirmed (Milvus accepts \texttt{ef=0} and \texttt{nprobe=0}; Qdrant accepts \texttt{m=0} and \texttt{bits=0}). We first asked DeepSeek to formalize each passage independently; it reproduced an over-strict clause on the same parameter in 6 of the 18 (the original five plus Qdrant \texttt{bits}; the other five scaling-pass clauses are family-specific)---the phenomenon concentrates in default-based parameters (\texttt{bits} default 8 joins the original default-based TI set), while example-based search parameters are family-specific. We then asked DeepSeek to judge each of GLM's clauses against the documentation; cross-model judging caught 8 of the 18 but missed 3 of the 6 task-intrinsic ones, where DeepSeek read the ambiguous documentation as GLM had. Source-grounded falsification contradicts all eighteen, because the implementation accepts the value each over-strict clause rejects. The TI rate is 6/18 (Wilson 95\% CI $[16\%, 56\%]$), tighter than the twelve-clause pilot's $[19\%, 68\%]$. This vendor-wise distribution is itself a finding: over-strict concentrates in APIs with many optional-default parameters (Milvus, Qdrant's search parameters) and is absent where documentation states explicit bounds (a parallel expansion on Weaviate v1.38.2 surfaces no over-strict clauses, because Weaviate documents ``Must be >= 1'' for \texttt{ef}, \texttt{dynamicEfMin}, etc., so its doc-code gaps are documentation-implementation bugs rather than over-formalized clauses), which is why the set is Milvus-heavy and why scaling beyond it is bounded by the phenomenon rather than by sampling effort. A within-vendor contrast isolates documentation style as the driver rather than the vendor: on Qdrant, the three search parameters that take optional defaults (\texttt{timeout}, \texttt{group\_size}, \texttt{score\_threshold}) are over-strict, while its collection parameters that state explicit minimums (\texttt{shard\_number}, \texttt{replication\_factor}, \texttt{write\_consistency\_factor}, all documented ``Minimum 1'') reject 0 and are not. The same contrast holds on Milvus, an independent vendor: its search and load parameters that take optional defaults (\texttt{ef}, \texttt{nprobe}, \texttt{level}, \texttt{replicaNumber}) are over-strict, while its parameters that state explicit ranges (\texttt{M}, enforced at $[1,2048]$; \texttt{nlist} at $[1,65536]$; \texttt{dimension} at $[1,32768]$) reject the violating value and are not. This yields a falsifiable prediction: a VDBMS parameter documented with an optional default and no explicit bound is a candidate for over-formalization, whereas one documented with an explicit minimum is not.

A behavior-level probe extends the phenomenon beyond parameters. On the original four Milvus by-design behavior issues (search on an unloaded collection, duplicate collection creation, drop on a non-existent collection, and a leading-underscore collection name, all maintainer-marked resolution/by-design) plus seven additional idempotency behaviors we probed---drop/create on a non-existent or already-existing partition, release on a not-loaded collection, create on an already-existing index, and load on an already-loaded collection (Milvus), Qdrant delete on a non-existent collection, and Weaviate delete on a non-existent class---each impl-confirmed to return success where the documentation implies an error, DeepSeek independently over-formalized the documented constraint on 11 of 11; cross-model judging caught only 1 of 11; source-grounded falsification caught all 11. The behavior-TI rate (11/11, Wilson 95\% CI $[74\%, 100\%]$) is the headline finding of this probe and is far higher than the parameter-TI rate (6/18) for a mechanistic reason: parameter documentation states a default (e.g., ``Default 1'') that only sometimes invites a bound, whereas behavior documentation states an explicit error condition (e.g., ``deleting a non-existent collection returns an error'') that both families uniformly over-formalize into a hard rejection, which the idempotent implementation then violates---so the two subtypes are reported separately, not pooled as the headline. The seven new behaviors are impl-confirmed idempotent but not separately maintainer-acknowledged, which we flag to preserve the within-vendor parameter contrast above. Notably, Weaviate contributes this idempotency-behavior phenomenon (delete on a non-existent class) while its parameters remain explicit-bound with zero over-strict, confirming both that the explicit-bound predictor holds cross-vendor and that the behavior phenomenon is not Milvus/Qdrant-specific. As an aggregate only, the pooled task-intrinsic rate across both subtypes ($n{=}29$) is $17/29$ (Wilson 95\% CI $[41\%, 75\%]$), and cross-model judging catches $4$ of the $17$.

A specificity check confirms the phenomenon does not generalize to explicit bounds. On twenty-one parameters whose documentation states explicit bounds---Qdrant's \texttt{shard\_number}, \texttt{replication\_factor}, \texttt{write\_consistency\_factor}, \texttt{full\_scan\_threshold}, \texttt{ef\_construct}, \texttt{indexing\_threshold}, \texttt{max\_optimization\_threads}, \texttt{wal\_segments\_ahead}, \texttt{max\_indexing\_threads}; Weaviate's \texttt{ef}, \texttt{dynamicEfMin}, \texttt{dynamicEfMax}, \texttt{efConstruction}, \texttt{maxConnections}, \texttt{vectorIndexType}, \texttt{replicationConfig.factor}; Milvus's \texttt{dimension}, \texttt{num\_partitions}, \texttt{M}, \texttt{efConstruction}, \texttt{nlist}---DeepSeek independently over-formalized zero ($0/21$, Wilson 95\% CI $[0\%, 16\%]$). The task-intrinsic phenomenon concentrates in ambiguous optional-default APIs and is absent where documentation states explicit bounds, reinforcing the within-vendor contrast.

\begin{table}[t]
\caption{Cross-model judging vs source-grounded falsification on eighteen GLM over-strict clauses (thirteen Milvus + five Qdrant v1.18.2, live-probe-confirmed). TI marks the six task-intrinsic clauses. Cross-model judging misses 3 of 6 TI clauses; source-grounded falsification contradicts all eighteen. The behavior over-strict (now eleven, all TI) and explicit-bound subtypes extend this probe to $n{=}50$ (see text, RQ3).}
\label{tab:e2}
\small
\begin{tabular}{@{}lp{0.10\columnwidth}p{0.20\columnwidth}p{0.18\columnwidth}@{}}
\toprule
Over-strict clause & TI & Cross-model judging & Source-grounded \\
\midrule
shardsNum $\geq 1$ & yes & missed & caught \\
metricType strict enum & no & missed & caught \\
consistencyLevel strict enum & no & caught & caught \\
data non-empty & yes & missed & caught \\
outputFields required & no & caught & caught \\
limit $\geq 1$ & no & caught & caught \\
roundDecimal $\geq 0$ & no & caught & caught \\
offset $\geq 1$ & no & caught & caught \\
dbName required & no & caught & caught \\
\midrule
\texttt{timeout} $\geq 1$ & yes & caught & caught \\
\texttt{group\_size} $\geq 1$ & yes & missed & caught \\
\texttt{score\_threshold} $\in[0,1]$ & yes & missed & caught \\
\midrule
\multicolumn{4}{l}{\emph{Scaling passes (contract-formalizer, live-probe-confirmed; 1/6 task-intrinsic):}} \\
\texttt{ef} $\geq 1$ & no & missed & caught \\
\texttt{nprobe} $\geq 1$ & no & missed & caught \\
\texttt{level} $\geq 0$ & no & caught & caught \\
\texttt{replicaNumber} $\geq 1$ & no & missed & caught \\
\texttt{m} $\geq 2$ & no & missed & caught \\
\texttt{bits} $\geq 1$ & yes & missed & caught \\
\midrule
Total & 6 & 8/18 & 18/18 \\
\bottomrule
\end{tabular}
\end{table}

\paragraph{Threats to validity.}
\emph{Internal validity.} The RQ3 probe covers fifty clauses across three subtypes (eighteen over-strict parameters, eleven by-design/idempotency behaviors, twenty-one explicit-bound negatives); DeepSeek over-formalizes at 17/29 on ambiguous optional-default APIs and 0/21 on explicit bounds. The over-strict subset remains the most contingent finding; the retrospective and the yield are the broader evidence base, and a larger head-to-head study is ongoing. \emph{External validity.} Generalization to Weaviate is breadth-only; statistical claims rest on Milvus and Qdrant, and result correctness of vector search (ANN recall, ranking) is out of scope. \emph{Construct validity.} All source-anchor results use a single model family (GLM-5.2); a cross-model check re-running the dev-reviewer with DeepSeek on twenty candidates (Cohen's $\kappa = 1.0$, blind to GLM-5.2's rationale, spanning input-validation, upsert-semantics, idempotent-drop, correct-reject, and dynamic-field subtypes) suggests the verdict is not strongly family-specific when source evidence is explicit; we still do not estimate recall because there is no public ground-truth defect catalog for VDBMSs.

\section{Related Work} \label{sec:related}

\paragraph{VDBMS testing.} VDBFuzz~\cite{vdbfuzz26} is the first dedicated VDBMS fuzzer and uses crash as its oracle, which makes it complementary to this work. The roadmap~\cite{roadmap25} sets the agenda we build on, and the empirical bug study~\cite{bugstudy25} the taxonomy; both point to the documentation-implementation subset we target.

\paragraph{REST-API oracle generation.} A recent line derives test oracles for REST APIs from structured sources: AGORA+~\cite{agoraplus25} from execution traces, SATORI~\cite{satori25} from OpenAPI specifications, and MASTOR~\cite{mastor26} from source code. As detailed in Section~\ref{sec:regime}, these operate in a low-ambiguity regime where the LLM transcribes explicit constraints rather than interpreting ambiguous documentation; MASTOR, the closest, tests implemented behavior rather than the documentation-implementation gap TestVDB targets.

\paragraph{LLM-as-judge reliability.} Self-preference bias, in which an LLM evaluator favors its own family's outputs, is established for general text evaluation (Panickssery et al.~\cite{panickssery24}; Wataoka and Takahashi~\cite{wataoka24}); we observe the analogous bias in the test-oracle pipeline. Our family-specific layer is an instance of it, and cross-model validation mitigates it as expected. A separate reliability axis is intra-judge self-inconsistency: Haldar et al.~\cite{haldar25} show that a single LLM judge's ratings vary substantially across identical runs. Our task-intrinsic layer is orthogonal to both: when the documentation is ambiguous, different model families converge on the same wrong clause at the extraction step, so the error survives cross-model validation. Section~\ref{sec:regime} scopes this extraction-level stability claim precisely.

\paragraph{Database and property-based oracles.} NoREC~\cite{norec20}, TLP~\cite{tlp20}, DQE~\cite{dqe20}, and DDLCheck~\cite{ddlcheck25} test SQL correctness against reference semantics that VDBMS documentation-implementation consistency lacks. Property-based testing~\cite{claessen00} and its REST instances, Schemathesis~\cite{schemathesis} and QuickREST~\cite{quickrest20}, together with metamorphic testing (MeTMaP~\cite{metmap24}), cover the mathematical-invariant subset; the documentation-implementation residual is where they, by construction, stop.

\paragraph{Documentation-derived oracles.} Toradocu~\cite{toradocu16} pioneered oracle extraction from natural-language documentation, using NLP and pattern matching to translate Javadoc \texttt{@throws} comments into executable assertions. Its extraction is deterministic but limited to simple syntactic patterns, and it acknowledges false positives from extraction failures without correcting them. Doc2OracLL~\cite{doc2oracll25} extended this line to LLMs, showing that documentation quality directly impacts LLM-generated oracle correctness, and AugmenTest~\cite{augmentest25} similarly infers oracles from the available documentation, but both still trust the generated oracle without verification. Konstantinou et al.~\cite{konstantinou24} provide direct evidence for the risk this creates: LLM-generated oracles tend to capture the actual program behavior rather than the expected, correct behavior, which is precisely the failure mode our source-grounded falsification targets. ChatAssert~\cite{chatassert24} addresses false positives through iterative prompt repair guided by compilation and execution feedback, and Testora~\cite{testora26} uses natural-language PR descriptions as a regression oracle, achieving 55\% precision even with a multi-question classifier, further evidence that LLM judgment alone remains unreliable. All of these treat the LLM as the final semantic arbiter: ChatAssert and Testora verify through runtime behavior (compilation feedback, differential execution), but the LLM still makes the final call. TestVDB differs by using the implementation's source code, not runtime behavior, as an independent verification source, falsifying LLM-derived behavioral claims and targeting task-intrinsic documentation-interpretation errors that prompt refinement and runtime repair cannot reach. Where these predecessors extract from low-ambiguity structured documentation (Javadoc \texttt{@throws} patterns, API specs), TestVDB targets the high-ambiguity prose regime whose task-intrinsic interpretation errors (Section~\ref{sec:regime}) they do not address.

\section{Discussion and Limitations} \label{sec:disc}

The setting is not specific to VDBMSs in principle, though our evaluation is VDBMS-only. Systems whose documentation is natural-language prose rather than a structured specification---REST APIs without OpenAPI coverage, configuration validation, and policy-as-code checks where the documentation is the only oracle---enter the same regime, where a single LLM family that both extracts the specification and checks documentation-implementation consistency is susceptible to the two-layer unreliability we describe, and source-grounded falsification is the natural mitigation. We claim transferability only to structurally similar documentation regimes and have not tested these transfers; empirical validation is future work.

Beyond the threats to validity in Section~\ref{sec:eval}, two limitations bound the approach itself. Source-grounded falsification requires source, so it does not transfer to closed-source VDBMSs, and it treats the implementation as correct, so an implementation bug can wrongly falsify a clause whose documentation is right. The association between documentation style (optional-default versus explicit-bound) and over-formalization risk is a correlative observation, validated by within-vendor contrast and a falsifiable prediction (RQ3) but not a causal claim; ruling out alternative explanations such as team structure or implementation complexity is future work. The 85\% residual is the composition of TestVDB's findings, biased toward documentation-implementation defects by design, not an estimate of the true defect distribution.

\paragraph{Crash-oracle asymmetry and incomplete fixes.} The two reachability directions are not symmetric. We hypothesize that the reason is structural: a crash on an input the documentation declares valid is itself a documentation-implementation violation, so a documentation-implementation oracle can reach crash-class defects at the subset where the crash-triggering input violates a documented bound, whereas the reverse need not hold, because a silent accept is not a crash and a crash oracle cannot fire on it. Our two controlled cases are consistent with this hypothesis but, at $n{=}1$ per direction, do not establish it. The temporal consequence, however, is visible in our data independent of that hypothesis. Qdrant issue \#9045 (TestVDB, acknowledged) root-causes a production panic (\#7967)\footnote{\url{https://github.com/qdrant/qdrant/issues/7967}} to an inconsistent validation path---\texttt{wait=true} rejects a zero-length vector while \texttt{wait=false} accepts it, because vector validation runs only under a \texttt{debug\_assert} that release builds skip---and prescribes an API-boundary fix. The panic symptom had recurred from v1.16.3 through v1.18.0 under crash-site patches that did not touch the validation gap, and was resolved only by the boundary fix. Crash-focused fuzzing surfaces symptoms whose fixes can leave the documentation-implementation residual intact, and TestVDB targets exactly that residual.

\section{Conclusion}

VDBMS API documentation-implementation consistency is a defect class where a substantial portion of accept/reject decisions resist deterministic checking, which leads us to adopt an LLM for extraction and judgment. The LLM's documentation-interpretation errors split into a family-specific layer that cross-model validation mitigates and a task-intrinsic layer that it cannot, because the ambiguity lives in the documentation. Source-grounded falsification resolves the task-intrinsic layer by treating the LLM-derived behavioral claims as hypotheses and the implementation's actual behavior as the reference. TestVDB demonstrates the approach at scale: 107 candidate issues across three VDBMSs, 49 true-positive defects (15 merged-PR-fixed, 16 with open fix-PRs, 18 acknowledged but unfixed), with the source-grounded dev-reviewer reaching 67\% precision and 74\% recall on Milvus and Qdrant. The boundary we draw, between problems where extraction from structured sources suffices and those where natural-language documentation forces LLM interpretation, is where we see the testing of LLM-dependent systems heading.

\bibliographystyle{ACM-Reference-Format}
\bibliography{references}

\end{document}
